\documentclass[11pt,a4paper]{article}
\usepackage[margin=25mm]{geometry}
\usepackage{amsmath,amssymb,amsthm}
\setlength{\emergencystretch}{2em}
\newtheorem{proposition}{Proposition}
\newtheorem{lemma}[proposition]{Lemma}
\newcommand{\gen}[1]{\langle#1\rangle}
\newcommand{\Alt}{\operatorname{Alt}}
\title{Opposite exceptional ends after a later unshared cut}
\author{Human Compression project: bounded L044 continuation}
\date{}
\begin{document}\maketitle
We retain the standing L044 hypotheses: $G\le A_n$, $n\ge12$,
is a transitive string C-group with connected diagram and every
maximal omission intransitive. The alternating rank bound at smaller
degrees at least $12$ is an explicit induction hypothesis.
Assume $r>\lfloor(n-1)/2\rfloor$ and an actual extendable
six-point source Table~2(2) or (3) head at the first three labels.
The accepted critical-head theorem gives $n=2r\ge18$.

\begin{proposition}[What remains at a later unshared cut]
Suppose $j\ge5$ has a unique eligible crossing, with actual supported
transitive sides $P=G_{<j}$ on $C$ and $Q=G_{>j}$ on $F$, the
original head contained in $C$. Then precisely the following
alternatives remain unexcluded here:
\begin{enumerate}
\item $j=r-1$, $F$ is a singleton, and $P=\Alt(C)=A_{n-1}$ naturally;
\item $j=r-3$, $Q$ is the degree-five $D_{10}$ marking, and
      $P=\Alt(C)=A_{n-5}$ naturally. The complete attachment theorem
      makes the actual last three marks a six-point Table~2 head.
\end{enumerate}
In particular an opposite actual six-point $A_5$ side is impossible.
These are residual alternatives, not assertions that either exists.
\end{proposition}
\begin{proof}
The large-pure-prefix theorem identifies $P=\Alt(C)$; write
$c=|C|$, $f=|F|$, and $d=\operatorname{rank}(Q)$.
It also supplies the endpoint anchors and local intransitive omissions.
Here $c\ge12$: for $8\le c<12$ we have $j\ge5\ge(c-1)/2$.
Apply the complete small primitive \emph{prefix exclusion}
(Theorem~3 of \emph{Small primitive attachments through a decorated
tail}), not the raw Table~2 list. Its hypotheses hold because $P$
is a proper actual prefix, the ambient degree is at least $12$,
and the ambient marking is even, connected and has intransitive
omissions. It excludes in particular the degree-ten rank-five and
degree-eleven rank-six rows, leaving only the degree-five or
degree-six exceptions, impossible on $C$. The small anchored
imprimitive bound covers a nonprimitive side. Thus induction gives
$j\le\lfloor(c-1)/2\rfloor$.

Except for a singleton, the full dual of that proper-prefix exclusion and
the proper primitive and anchored imprimitive bounds give
$d\le\lfloor(f-2)/2\rfloor$, unless $Q$ is a degree-five $D_{10}$,
a degree-six $A_5$ head, or a full natural alternating group of
degree at least $12$. The strict inequality case gives
$r=j+d+1\le(n-1)/2$, contrary to the assumption.
When both sides are large natural alternating groups, the complete
faithful even one-point extension lemma gives an actual side marking
of degree $c+1$ and rank $j+1$, or its counterpart on $F$.
Its complete smaller-degree bound, added to the other side's bound,
again gives $r\le\lfloor(n-1)/2\rfloor$. This is the exact paid
two-natural-side argument, including its projection-kernel proof.

For an opposite $A_5$ side, reverse the entire string. Its first
two marks generate the degree-five dihedral group with trivial
centralizer on that orbit. The crossing therefore fixes those five
points and has the form $(v\ b)\delta$. Its actual opposite
suffix is $P=\Alt(C)$. The accepted pure-head/natural-suffix
obstruction applies and excludes it. The impossible reversed
row~2 orientation is excluded by the literal head-extension lemma;
row~3 has its stated reversing conjugation.

Only the singleton and the $D_{10}$ alternative remain.
Their ranks are respectively zero and two, so $j=r-1$ or $r-3$.
In the latter case the complete six-choice dihedral attachment,
applied to the globally reversed string, gives the stated actual
last three marks. It does not assert that the next cut is unshared.
\end{proof}

\begin{lemma}[The actual middle between two $A_5$ ends]
Suppose the actual first and last three marks are extendable
six-point Table~2 heads. Let their supports be $D_L,D_R$, and let
$v_L,v_R$ be their exit vertices. The supports are disjoint, and
\[
 G_{\{3,\ldots,r-4\}}=\Alt(Z),\qquad
 Z=\Omega\setminus((D_L\setminus\{v_L\})
                   \cup(D_R\setminus\{v_R\})).
\]
This is an actual natural alternating group of degree $n-10$
and rank $r-6=(n-12)/2$. It has both endpoint anchors and all
maximal omissions intransitive on $Z$.
\end{lemma}
\begin{proof}
Since $r\ge9$, the two heads commute. A permutation centralizing
either transitive six-point head preserves its unique nontrivial
orbit and fixes that orbit pointwise, by its trivial centralizer.
Thus the heads have disjoint supports.
The complete critical-head theorem and its full dual give
\[
 G_{\ge3}=\Alt(\Omega\setminus(D_L\setminus\{v_L\})),\qquad
 G_{\le r-4}=\Alt(\Omega\setminus(D_R\setminus\{v_R\})).
\]
Their natural permutation-group intersection is $\Alt(Z)$;
full ambient IP identifies it with the displayed middle interval.
No quotient IP is used. Every middle mark except its first fixes
$v_L$, and every one except its last fixes $v_R$.
If one of its maximal omissions were transitive on $Z$, adjoining
the two full heads would make the corresponding ambient omission
transitive on all of $\Omega$, a contradiction. The rank count
is literal. For even degree $n-10\ge12$ it is equality in, not
a contradiction to, the smaller-degree alternating rank bound.
\end{proof}

\begin{lemma}[A bounded exclusion within the two-ended alternative]
The preceding two-ended configuration requires $n\ge24$.
\end{lemma}
\begin{proof}
Use the left critical-head notation
\[
 \tau=\rho_3=(v_L\ b)\delta,\quad \sigma=\rho_4,
 \qquad G_{\ge3}=\Alt((\Omega\setminus D_L)\cup\{v_L\}).
\]
The suffix $G_{\ge5}$ fixes $v_L,b$; transitivity of this critical
tail forces $x=b^\sigma\ne b$. The actual pure-head overlap proof
gives $\Alt(W_L)\le G_{\le4}$, where $W_L=D_L\cup\{b,x\}$.
It also forces $\tau x=x$: otherwise a second overlap supplies a
pure $A_4$ inside
$G_{\le4}\cap G_{\ge3}=\gen{\tau,\sigma}$, a dihedral group,
which cannot contain noncommuting three-cycles.
Hence $\sigma=(b\ x)\epsilon$, where $\delta,\epsilon$ are odd
involutions fixing $W_L$ pointwise. The same actual intersection
contains a pure three-cycle on $\{v_L,b,x\}$.
Writing $m=|\delta\epsilon|$, the disjoint core/exterior product
of $\tau\sigma$ shows that $3\nmid m$: a pure core rotation
requires a multiple of $m$ not divisible by $3$.
If $m\le2$, then $\delta,\epsilon$ commute, so $\delta$
centralizes $\tau,\sigma,G_{\ge5}$ and fixes $v_L$, contrary to
transitivity of the critical tail and oddness. Thus $m\ge4$.
In particular $J_L=\gen{\delta,\epsilon}$ is nonabelian.

Both involutions also fix $D_R$ pointwise. On at most five points
an odd involution is a transposition, and two transpositions have
product order at most three. Their exterior support therefore
needs at least six points. If $n=18$, only $n-14=4$ points are
available outside $W_L\cup D_R$, an immediate contradiction.
If $n=20$, the right crossing $\rho_{r-4}=\rho_6$ commutes
with all of $G_{\le4}$ and hence fixes $W_L$ pointwise.
It sends $v_R$ to a point outside $D_R\cup W_L$ which is fixed
by $\tau$ and $\sigma$, because both fix $v_R$.
That point is also fixed by $\delta,\epsilon$, leaving at most
five points for their nontrivial action, again impossible.

Finally let $n=22$, so $r=11$. The right-hand version supplies
$W_R$ of size eight and a nonabelian group $J_R$ generated by
two odd exterior involutions. The entire prefixes $G_{\le4}$
and $G_{\ge6}$ commute, so $W_L,W_R$ are disjoint, and both
exterior groups act on the remaining six-point set $V$.
Their generators commute, since the original marks commute and
the removed core transpositions have disjoint supports.

On six points, an odd involution has either one or three
transpositions. Two single transpositions have product order at
most three. Two fixed-point-free involutions also have product
order at most three, as their union is a disjoint union of
alternating even cycles of length at most six.
Thus one of $\delta,\epsilon$ is a transposition and the other
a perfect matching. If the former is an edge of the latter,
their product has order at most two. Otherwise their joint orbits
have sizes four and two, with the standard $D_8$ action on four
points and $C_2$ on two points. Its centralizer in $S_V$ lies in
$C_2\times C_2$: on the four-point orbit, centralizing a
four-cycle and a reflection leaves only its central involution;
the two orbit sizes prevent an exchange. This centralizer is
abelian, contradicting the nonabelian $J_R$ centralizing $J_L$.
All possibilities $18\le n<24$ are excluded.
\end{proof}

\paragraph{Exact source route and remaining work.}
This refines the exceptional continuation of source
Propositions~5.28 and 5.32 (physical pages 36--38 and 44--45).
The full supported-side hypotheses remain explicit.
The $D_{10}$ alternative creates an actual opposite $A_5$ head,
not an automatically unshared next crossing. The resulting
natural middle and its anchors are proved, but its maximal
permitted rank is not asserted to classify its endpoint markings.
The two-ended case for $n\ge24$, the singleton alternative and
shared-label crossings remain open here. Existing complete
small-degree certificates and all full predecessor proofs stay
paid; no new finite computation or formalization is a premise.
No whole L044, original or question credit is claimed.
\end{document}
